\documentclass[10pt]{article}

\usepackage[utf8]{inputenc}

\usepackage[T1]{fontenc}

\usepackage{amsmath}

\usepackage{amsfonts}

\usepackage{amssymb}

\usepackage[version=4]{mhchem}

\usepackage{stmaryrd}

\usepackage{bbold}



\begin{document}

в области, называемой областью зависимости отрезка $M N$. Эти же начальные значения влияют на решение в других областях, хотя и не полностью определяют его. Эти области называются областями влияния отрезка $M N$; для остальных областей никак не влияют начальные условия на отрезке $M N$.



Если же $B^{2}-A C<0$, то характеристики отсутствуют. При этом уравнение (1) относят к эллиптич. типу. В этом случае решение $u(x, t)$ уравнения (1) непрерывно вместе со всеми своими производными во всей области существования.



Jum.: [1] К у р а н т Р., Уравнения с частными производными, пер. с англ., М., 1964; [2] Т ихо но в А. Н., С а м а р ск и й А. А., Уравнения математической физики, 5 изд., М., 1977; [3] Го д у но в С. К., Уравнения математической физики, 2 изд., М., 1979; [4] В ла д и м и р о в В. С., Уравнения математической физики, 4 изд., М., 1981. В. П. Демуцкий, Р. В. Половин.



XAPAKTEPИСTИКА гипе рповерхности в кон тактном пространстве с контактной формой $\alpha-$ интегральная кривая поля ядер ограничения дифференциала $d \alpha$ на пересечение гиперплоскости контактного поля с касательным пространством гиперповерхности. Характеристич. поле не зависит от выбора контактной формы и имеет особенности в точках касания гиперповерхности контактного поля.



М. Э. Казарян.



XАPAKTEPИСTИКА гипе рпове рхности в симплект и ческом п р о с т р а с т ве - интегральная кривая поля ядер ограничения симплектической структуры на эту гиперповерхность. Если гиперповерхность является поверхностью уровня гамильтониана $H$, то характеристич. поле направлений задается ограничением на гиперповерхность гамильтонова векторного поля с гамильтонианом $H$.



М. Э. Казарян.



XАPAKTEPИCTИKА д и ф ф е р е н ц и а л ь н о го о п е р а то ра - понятие теории дифференциальных уравнений с частными производными. Пусть $\hat{H}-\hat{A}$-псевдодифференциальный оператор с главным символом $H(x, p)$. Х а ра к теристиками оператора $\hat{H}$ называются поверхности $S(x)=$ const с функцией $S$, удовлетворяющей уравнению Гамильтона - Якоби $H\left(x, \frac{\partial S}{\partial x}\right)=0$. Б и х а р а к т е р и с т и кам и этого оператора называются траектории системы Гамильтона для гамильтониана $H(x, p)$ в фазовом пространстве:



\[

\dot{x}=H_{p}(x, p), \dot{p}=-H_{x}(x, p) .

\]



Лum.: [1] М а с ло в В. П., «Успехи математич. наук», 1983, т. 38, B. 6 , c. 3-36.



Б. Ю. Стернин, В. Е. Шаталов.



ХАРАКТЕРИСТИКА контактн о Й формы - кривая в контактном многообразии, касательная прямая которой в каждой точке является ядром 2-формы $d \lambda$, где $\lambda$ - контактная форма. X. трансверсальны контактной структуре. Векторное поле $v$, касающееся Х., называется х а а кте ристическим полем (полем Риба), если $\lambda(v)=1$. Такое поле сохраняет контактную структуру. Пример X.геодезич. поток в многообразии контактных элементов риманова многообразия.



2-форма $d \lambda$ определяет в многообразии (определенном лишь локально) Х. симплектич. структуру. С. А. Табачников.



XАРАКТEPИСТИКА уравнения с частными произ водны м и $F(u, \partial u / \partial x, x)=0, x \in \mathbb{R}^{n}-$ характеристика гиперповерхности в контактном пространстве 1-струй функций на $\mathbb{R}^{n}$, заданной уравнением $F(u, p, x)=0$. Смысл $X$. состоит в том, что 1-струя любого решения целиком состоит из X. В координатах Х. задаются дифференциальными уравнениями



\[

\dot{x}=F_{p}, \dot{p}=-F_{x}-p F_{u}, \dot{u}=p F_{p} .

\]



См. также Характеристик метод.



М. Э. Казарян.



ХАРАКТЕРИСТИЧЕСКАЯ ФУНКЦИЯ с Л У ч а й н о Й в е л и ч и ны - Фурье преобразование функции распределения этой случайной величины; удобный аналитический объект для математических исследований различных свойств случайной величины.



Пусть $\xi-$ случайная величина,



\[

F(x)=\mathrm{P}\{\xi<x\}

\]



\begin{itemize}

  \item ее функция распределения. Функция

\end{itemize}



\[

f(t)=f_{\xi}(t)=\langle\exp i t \xi\rangle=\int_{-\infty}^{\infty} \exp (i t \xi) d F

\]



$(t \in \mathbb{R})$ называется х а акте р истиче ской функцие й случайной величины $\xi$. В случае, когда случайная величина обладает плотностью распределения вероятностей $p(x)$, интеграл (1) принимает вид



\[

f(t)=\int_{-\infty}^{\infty} \exp (\text { itx }) p(x) d x .

\]



Х. ф. $f(t)$ любой случайной величины обладает следующими свойствами:



$f(0)=1$,



$f(t)$ - непрерывна,



$f(t)$ - положительно определена, то есть



\[

\sum_{i, j} f\left(t_{i}-t_{j}\right) z_{i} z_{j}^{*} \geqslant 0

\]



где $t_{1}, \ldots, t_{s}-$ произвольный конечный набор значений аргумента $t$, а $z_{1}, \ldots, z_{s}$ - произвольный набор комплексных чисел $\left({ }^{*}\right.$ - означает комплексное сопряжение). Указанные три свойства являются определяюшими, а именно, для любой функции $f(t)$, обладающей этими тремя свойствами, найдется (и при этом лишь одна) функция $F(x)$ распределения вероятностей значений нек-рой случайной величины $\xi$ такая, что $f(t)$ представима в виде (1).



Справедливы равенства для моментов



\[

m_{n} \equiv \mathrm{E} \xi^{n}=(-i)^{n} f^{(n)}(0),

\]



где $f^{(n)}(0)$ - производная функции $f(t)$ порядка $n$.



В случае, когда у случайной величины существует $k$ моментов $m_{n}, n=1, \ldots, k$, функция $\ln f(t)$, по крайней мере, $k$ раз дифференцируема в малой окрестности точки $t=0$; ее производные в этой точке:



\[

\left.\frac{d^{k} \ln f(t)}{d t^{k}}\right|_{t=0},

\]



называются семи ин вариантам и (кумулянтам и) случайной величины.



Для двух независимых случайных величин $\xi_{1}$ и $\xi_{2}$ с Х. Ф. $f_{1}(t)$ и $f_{2}(t)$, Х. ф. их суммы $\eta=\xi_{1}+\xi_{2}$



\[

f_{\eta}(t)=f_{1}(t) f_{2}(t), t \in \mathbb{R} .

\]



Это свойство чрезвычайно существенно в теории суммирования случайных величин. Из него следует, что при добавлении к случайной величине константы $b$ Х. Ф. случайной величины умножается на $\exp ($ ita $)$,



\[

f_{\xi+b}(t)=\exp (i t b) f_{\xi}(t) .

\]



Для наиболее употребительных случаев гауссовского, пуассоновского и равномерного распределений их Х. Ф. вычислены и имеют вид:



для Гаусса распределения со средним $a$ и дисперсией $\sigma^{2}$ X. ф.



\[

f(t)=\exp \left(i t a-\sigma t^{2}\right)

\]



для Пуассона распределения с интенсивностью $\lambda$



\[

f(t)=\exp \{\lambda[\exp i t-1]\}

\]



для равномерного распределения на конечном отрезке $[a, b] \subset \mathbb{R}\left[\xi \in[a, b]\right.$, плотность $\left.p(x)=(b-a)^{-1}\right]$



\[

f(t)=\frac{\exp i t b-\exp i t a}{i t(b-a)} .

\]



В случае нескольких случайных величин $\xi_{1}, \ldots, \xi_{s}$ их совме-





\end{document}